% GENERATED FILE. Edit canonical lessons, then run npm run generate:books.
% canonical-source-hash: fnv1a64:b0fd6675f4060922
% canonical-lessons: MR-C254-corne, MR-C254-sodhun-kadhne, MR-C254-seva-karne, MR-C254-sahan-karne, MR-C254-svapna-pahne

\chapter{To steal, To discover, To serve, To bear, To dream}
\label{ch:mr-to-steal-to-discover-to-serve-to-bear-to-dream}
\hlchaptermodality{254}

\begin{chapteropening}
\textbf{By the end of this chapter:} \emph{I can say to steal, to discover, to serve, to bear and to dream.}

Complete the last lesson of chapter 254: I can say to steal, to discover, to serve, to bear and to dream.
\end{chapteropening}

\section[\texorpdfstring{\mr{चोरणे} (corṇe)}{चोरणे (corṇe)}]{\mr{चोरणे} (corṇe) — to steal}

\label{lesson:MR-C254-corne}



\begin{quote}
Before the new one: say the Marathi for to inform, then the Marathi for to complain.
\end{quote}

\subsection*{You'll want to know: \mr{चोरणे}}
\textbf{\mr{चोरणे}} — \emph{corṇe} — \textquotedblleft{}to steal\textquotedblright{}.

\begin{grammarlens}[title={What you've built}]
One more word for this chapter.
\end{grammarlens}

\subsection*{Your turn}
\begin{itemize}
\raggedright
  \item \emph{Say it:} \emph{corṇe}
  \item \emph{Say it:} \emph{corṇe}, once more
  \item \emph{Say it:} say \emph{corṇe}
  \item \emph{Recall:} say the Marathi for to inform, then the Marathi for to complain, then say \emph{corṇe} again
\end{itemize}

\begin{tcolorbox}[breakable,colback=teal!4,colframe=teal!35!black,title={Before you move on}]
What does \mr{चोरणे} mean? (\textquotedblleft{}To steal\textquotedblright{}.) Say it once more.
\end{tcolorbox}
\section[\texorpdfstring{\mr{शोधून} \mr{काढणे} (śodhūn kāḍhṇe)}{शोधून काढणे (śodhūn kāḍhṇe)}]{\mr{शोधून} \mr{काढणे} (śodhūn kāḍhṇe) — to discover}

\label{lesson:MR-C254-sodhun-kadhne}



\begin{quote}
Before the new one: say the Marathi for to complain, then the Marathi for to steal.
\end{quote}

\subsection*{You'll want to know: \mr{शोधून} \mr{काढणे}}
\textbf{\mr{शोधून} \mr{काढणे}} — \emph{śodhūn kāḍhṇe} — \textquotedblleft{}to discover\textquotedblright{}.

\begin{grammarlens}[title={What you've built}]
One more word for this chapter.
\end{grammarlens}

\subsection*{Your turn}
\begin{itemize}
\raggedright
  \item \emph{Say it:} \emph{śodhūn kāḍhṇe}
  \item \emph{Say it:} \emph{śodhūn kāḍhṇe}, once more
  \item \emph{Say it:} say \emph{śodhūn kāḍhṇe}
  \item \emph{Recall:} say the Marathi for to complain, then the Marathi for to steal, then say \emph{śodhūn kāḍhṇe} again
\end{itemize}

\begin{tcolorbox}[breakable,colback=teal!4,colframe=teal!35!black,title={Before you move on}]
What does \mr{शोधून} \mr{काढणे} mean? (\textquotedblleft{}To discover\textquotedblright{}.) Say it once more.
\end{tcolorbox}
\section[\texorpdfstring{\mr{सेवा} \mr{करणे} (sevā karṇe)}{सेवा करणे (sevā karṇe)}]{\mr{सेवा} \mr{करणे} (sevā karṇe) — to serve}

\label{lesson:MR-C254-seva-karne}



\begin{quote}
Before the new one: say the Marathi for to steal, then the Marathi for to discover.
\end{quote}

\subsection*{You'll want to know: \mr{सेवा} \mr{करणे}}
\textbf{\mr{सेवा} \mr{करणे}} — \emph{sevā karṇe} — \textquotedblleft{}to serve\textquotedblright{}.

\begin{grammarlens}[title={What you've built}]
One more word for this chapter.
\end{grammarlens}

\subsection*{Your turn}
\begin{itemize}
\raggedright
  \item \emph{Say it:} \emph{sevā karṇe}
  \item \emph{Say it:} \emph{sevā karṇe}, once more
  \item \emph{Say it:} say \emph{sevā karṇe}
  \item \emph{Recall:} say the Marathi for to steal, then the Marathi for to discover, then say \emph{sevā karṇe} again
\end{itemize}

\begin{tcolorbox}[breakable,colback=teal!4,colframe=teal!35!black,title={Before you move on}]
What does \mr{सेवा} \mr{करणे} mean? (\textquotedblleft{}To serve\textquotedblright{}.) Say it once more.
\end{tcolorbox}
\section[\texorpdfstring{\mr{सहन} \mr{करणे} (sahan karṇe)}{सहन करणे (sahan karṇe)}]{\mr{सहन} \mr{करणे} (sahan karṇe) — to bear, to endure}

\label{lesson:MR-C254-sahan-karne}



\begin{quote}
Before the new one: say the Marathi for to discover, then the Marathi for to serve.
\end{quote}

\subsection*{You'll want to know: \mr{सहन} \mr{करणे}}
\textbf{\mr{सहन} \mr{करणे}} — \emph{sahan karṇe} — \textquotedblleft{}to bear, to endure\textquotedblright{}.

\begin{grammarlens}[title={What you've built}]
One more word for this chapter.
\end{grammarlens}

\subsection*{Your turn}
\begin{itemize}
\raggedright
  \item \emph{Say it:} \emph{sahan karṇe}
  \item \emph{Say it:} \emph{sahan karṇe}, once more
  \item \emph{Say it:} say \emph{sahan karṇe}
  \item \emph{Recall:} say the Marathi for to discover, then the Marathi for to serve, then say \emph{sahan karṇe} again
\end{itemize}

\begin{tcolorbox}[breakable,colback=teal!4,colframe=teal!35!black,title={Before you move on}]
What does \mr{सहन} \mr{करणे} mean? (\textquotedblleft{}To bear, to endure\textquotedblright{}.) Say it once more.
\end{tcolorbox}
\section[\texorpdfstring{\mr{स्वप्न} \mr{पाहणे} (svapna pāhṇe)}{स्वप्न पाहणे (svapna pāhṇe)}]{\mr{स्वप्न} \mr{पाहणे} (svapna pāhṇe) — to dream}

\label{lesson:MR-C254-svapna-pahne}



\begin{quote}
Before the new one: say the Marathi for to serve, then the Marathi for to bear.
\end{quote}

\subsection*{You'll want to know: \mr{स्वप्न} \mr{पाहणे}}
\textbf{\mr{स्वप्न} \mr{पाहणे}} — \emph{svapna pāhṇe} — \textquotedblleft{}to dream\textquotedblright{}.

\begin{grammarlens}[title={What you've built}]
One more word for this chapter.
\end{grammarlens}

\subsection*{Your turn}
\begin{itemize}
\raggedright
  \item \emph{Say it:} \emph{svapna pāhṇe}
  \item \emph{Say it:} \emph{svapna pāhṇe}, once more
  \item \emph{Say it:} say \emph{svapna pāhṇe}
  \item \emph{Recall:} say the Marathi for to serve, then the Marathi for to bear, then say \emph{svapna pāhṇe} again
\end{itemize}

\begin{tcolorbox}[breakable,colback=teal!4,colframe=teal!35!black,title={Before you move on}]
What does \mr{स्वप्न} \mr{पाहणे} mean? (\textquotedblleft{}To dream\textquotedblright{}.) Say it once more.
\end{tcolorbox}
